\section{Análisis ético y de sesgos}\label{sec:etica-sesgos}

Este apartado convierte la evidencia de calidad y representatividad del análisis exploratorio en controles para un eventual piloto. No entrena modelos, no estima desempeño ni atribuye discriminación algorítmica: esos resultados no existen todavía. La orientación es coherente con la necesidad de autonomía humana, transparencia, responsabilidad, inclusión y evaluación continua en sistemas de IA para salud \parencite{who-ai-health,fda-transparency-mlmd}.

\subsection{Evidencia, interpretación y límites}

La fuente original contiene 1.025 filas, de las cuales 723 son duplicados exactos (70,54\%); la versión analítica deduplicada conserva 302 registros. Esto constituye un problema de calidad y de trazabilidad del muestreo: la coincidencia exacta no permite conocer su causa. Si más adelante se entrenara o evaluara un modelo sin controlar este punto, las observaciones repetidas podrían sobreponderar patrones e incluso contaminar una separación de entrenamiento y prueba. No se afirma que dicha contaminación haya ocurrido. La mitigación exigible es conservar el original, declarar el criterio de deduplicación y verificar que registros equivalentes no queden repartidos entre particiones; el criterio CA-E01 formaliza su aceptación.

En los 302 registros únicos hay 96 mujeres (\texttt{sex}=0; 31,79\%) y 206 hombres (\texttt{sex}=1; 68,21\%). La codificación está respaldada por la fuente UCI \parencite{uci-heart-disease}. Esta diferencia describe cobertura desigual, no discriminación ni desempeño diferente. Su impacto técnico posible es que un modelo futuro disponga de menor evidencia para el grupo menos representado; por ello el eventual desempeño debe evaluarse por sexo con tamaño muestral suficiente y no se debe aprobar uso donde la evidencia no respalde la población objetivo.

La cobertura también es reducida en edades extremas: $<40$ aporta 15 registros únicos (4,97\%), 50--59 concentra 125 (41,39\%) y 70+ sólo 10 (3,31\%). Las intersecciones mujeres $\times <40$ y mujeres $\times 70+$ contienen cinco registros cada una. Estos tamaños hacen inestables los porcentajes descriptivos de \texttt{target}; no revelan prevalencia poblacional, causalidad ni discriminación. Las diferencias de \texttt{target} por sexo y edad ya registradas en \texttt{representatividad\_subgrupos.csv} deben leerse sólo con esa cautela y con el N correspondiente.

\subsection{Impacto ético y técnico del uso futuro}

El dataset de Kaggle sirve como evidencia académica inicial, pero no valida por sí mismo un uso en un hospital chileno. Pueden cambiar población, práctica clínica, medición y prevalencia; por ello la transferencia directa sería un riesgo de cambio de dominio. Cualquier paso hacia uso local requiere validación con datos y flujo clínico locales, además de monitoreo de deriva. Las guías para tecnologías clínicas de aprendizaje automático recomiendan comunicar límites de datos y poblaciones poco representadas, realizar validación específica del sitio y monitoreo continuo \parencite{fda-transparency-mlmd}.

Un resultado automatizado no puede operar como diagnóstico ni decisión autónoma. Falsos negativos podrían omitir una revisión preventiva y falsos positivos podrían generar alertas o carga innecesarias; ambos efectos requieren un profesional que interprete el contexto, pueda discrepar y deje trazabilidad de la decisión. El esquema \emph{Human in the Loop} es, por tanto, un requisito de control y no una característica opcional. SHAP puede apoyar explicaciones locales o globales, pero no prueba equidad ni reemplaza métricas por subgrupo, validación externa o supervisión clínica.

Privacidad y gobierno son igualmente éticos: antes de cualquier dato clínico deben definirse propósito, minimización, pseudonimización, control de acceso de mínimo privilegio, retención, auditoría y responsables. Cloud, Local e Híbrida presentan superficies operacionales distintas: la nube puede añadir transferencia y dependencia de proveedor; lo local concentra el control y la carga de operación; lo híbrido agrega sincronización y doble administración. Ninguna alternativa es intrínsecamente más ética; la aceptación depende de controles verificables y de la autorización que corresponda.

\subsection{Matriz de controles y criterios de aceptación}

La \tabref{tab:matriz-etica} sintetiza la cadena evidencia--impacto--mitigación--aceptación. La versión completa y reutilizable se encuentra en \texttt{outputs/tables/matriz\_etica.csv}; su generación desde los archivos original y deduplicado es reproducible mediante el script documentado en \texttt{src/}.

\begin{table}[H]
  \centering
  \caption{Matriz ética consolidada}
  \label{tab:matriz-etica}
  \scriptsize
  \resizebox{\textwidth}{!}{%
  \begin{tabular}{llll}
    \toprule
    ID & Hallazgo y evidencia & Impactos potenciales & Mitigación y aceptación \\
    \midrule
    E01 & \parbox[t]{0.27\textwidth}{Duplicación masiva y trazabilidad del tratamiento de datos. 1025 filas originales; 723 duplicadas exactas (70,54 \%); 302 únicas.} & \parbox[t]{0.29\textwidth}{Técnico: Sobreponderación de patrones y posible contaminación entre particiones si se modela sin control. Ético: Decisiones basadas en evidencia de calidad incierta.} & \parbox[t]{0.30\textwidth}{Conservar el original, documentar deduplicación y verificar separación por registro único antes de modelar. CA-E01: tratamiento de duplicados documentado y verificable.} \\
    E02 & \parbox[t]{0.27\textwidth}{Cobertura desigual por sexo. Únicos: mujeres (sex=0) 96 (31,79 \%); hombres (sex=1) 206 (68,21 \%).} & \parbox[t]{0.29\textwidth}{Técnico: Menor evidencia para estimar desempeño futuro en el grupo menos representado. Ético: Riesgo de uso inequitativo si no se evalúa la cobertura y el desempeño posterior.} & \parbox[t]{0.30\textwidth}{Reportar métricas por sexo sólo al existir modelo y tamaño suficiente; no extrapolar cobertura a población. CA-E02 y CA-E04.} \\
    E03 & \parbox[t]{0.27\textwidth}{Cobertura limitada en extremos etarios. Únicos: <40=15 (4,97 \%); 50--59=125 (41,39 \%); 70+=10 (3,31 \%).} & \parbox[t]{0.29\textwidth}{Técnico: Estimaciones inestables y generalización incierta en grupos con N reducido. Ético: Posible inequidad si se usa fuera de los grupos respaldados.} & \parbox[t]{0.30\textwidth}{Definir mínimo N, evaluar por edad cuando corresponda y restringir uso no respaldado. CA-E03 y CA-E04.} \\
    E04 & \parbox[t]{0.27\textwidth}{Evidencia interseccional insuficiente. Mujeres <40=5; mujeres 70+=5 registros únicos.} & \parbox[t]{0.29\textwidth}{Técnico: Porcentajes y métricas por intersección serían altamente inestables. Ético: Conclusiones aparentes podrían invisibilizar incertidumbre de grupos pequeños.} & \parbox[t]{0.30\textwidth}{Informar N junto con porcentajes; no aprobar uso en combinaciones sin evidencia suficiente. CA-E03 y CA-E04.} \\
    E05 & \parbox[t]{0.27\textwidth}{Diferencias descriptivas de target entre subgrupos. Las proporciones de target por sexo y edad están exportadas con N en representatividad\_subgrupos.csv.} & \parbox[t]{0.29\textwidth}{Técnico: No permiten inferir desempeño, prevalencia ni causalidad. Ético: Interpretación indebida como diferencia poblacional o discriminación.} & \parbox[t]{0.30\textwidth}{Tratar resultados como descriptivos; contrastar con validación clínica/local posterior. CA-E05: validación con datos locales antes de uso.} \\
    \bottomrule
  \end{tabular}%
  }
\end{table}

\begin{table}[H]
  \centering
  \caption{Matriz ética consolidada (continuación)}
  \label{tab:matriz-etica-cont}
  \scriptsize
  \resizebox{\textwidth}{!}{%
  \begin{tabular}{llll}
    \toprule
    ID & Hallazgo y evidencia & Impactos potenciales & Mitigación y aceptación \\
    \midrule
    E06 & \parbox[t]{0.27\textwidth}{Cambio de dominio Kaggle hacia hospital chileno. Dataset académico externo; no representa por sí mismo la población ni el flujo clínico local.} & \parbox[t]{0.29\textwidth}{Técnico: Deriva de población, medición o práctica clínica; desempeño no transferible. Ético: Riesgo de daños o acceso inequitativo por generalización injustificada.} & \parbox[t]{0.30\textwidth}{Validación local prospectiva, documentación de población objetivo y monitoreo de deriva. CA-E05 y CA-E09.} \\
    E07 & \parbox[t]{0.27\textwidth}{Automatización y posibles falsos negativos o positivos. El alcance define apoyo preventivo y revisión humana; no existe modelo ni métrica clínica validada.} & \parbox[t]{0.29\textwidth}{Técnico: Alertas erróneas u omisiones si un resultado se interpreta como decisión clínica. Ético: Daño potencial, sesgo de automatización y dilución de responsabilidad.} & \parbox[t]{0.30\textwidth}{Human in the Loop, protocolo de revisión, registro de decisión y uso exclusivamente preventivo. CA-E06 y CA-E10.} \\
    E08 & \parbox[t]{0.27\textwidth}{Privacidad y gobierno del dato sensible. Las alternativas Cloud, Local e Híbrida procesan datos de salud con exposiciones operacionales distintas.} & \parbox[t]{0.29\textwidth}{Técnico: Acceso indebido, retención excesiva o trazabilidad insuficiente. Ético: Afectación de confidencialidad y autonomía de las personas.} & \parbox[t]{0.30\textwidth}{Minimización, pseudonimización, mínimos privilegios, propósito definido, retención y auditoría. CA-E08.} \\
    E09 & \parbox[t]{0.27\textwidth}{Transparencia y explicabilidad insuficientes. SHAP está contemplado como apoyo explicativo; no hay modelo entrenado ni auditoría de desempeño.} & \parbox[t]{0.29\textwidth}{Técnico: No detectar límites o errores por subgrupo mediante explicaciones aisladas. Ético: Confianza injustificada y falta de rendición de cuentas.} & \parbox[t]{0.30\textwidth}{Documentar datos, límites, versiones y auditorías; SHAP complementa, no sustituye, pruebas de equidad. CA-E07 y CA-E09.} \\
    \bottomrule
  \end{tabular}%
  }
\end{table}


\subsection{Gates propuestos antes de un piloto}

Los criterios de aceptación no fijan umbrales inventados ni autorizan un modelo inexistente. Se proponen como condiciones de avance: CA-E01, tratamiento documentado de duplicados; CA-E02, desempeño por sexo cuando exista N suficiente; CA-E03, desempeño por rango etario cuando exista N suficiente; CA-E04, no aprobar uso en grupos sin cobertura respaldada; CA-E05, validación con datos locales; CA-E06, supervisión humana obligatoria; CA-E07, explicabilidad documentada; CA-E08, pseudonimización y control de acceso; CA-E09, monitoreo de deriva; y CA-E10, alcance limitado a apoyo preventivo, sin diagnóstico ni decisión automática. La aceptación final debe ser resuelta por la contraparte clínica y el gobierno del proyecto, con evidencia posterior.
